% ~10 min
\partframe{Machine-Checked}{Hylomorphisms in Rocq}

\begin{frame}
  \frametitle{What You Have to Trust}
  \begin{center}
    \begin{tikzpicture}[alt={The pipeline: a specification, written and
        read by a human, goes through calculation steps checked by the Rocq
        kernel, then through extraction to OCaml, then through the OCaml
        compiler. The specification, the kernel, extraction and the
        compiler are trusted. The calculation steps, and whoever chose
        them, are not.},
      stage/.style = {minimum height = 11mm, text width = 23mm,
        align = center}]
      \node[box, stage] (s) {specification};
      \node[box, stage, right = 9mm of s, draw = good, fill = good!8] (c)
        {calculation steps};
      \node[box, stage, right = 9mm of c] (e) {extraction to OCaml};
      \node[box, stage, right = 9mm of e] (o) {OCaml compiler};
      \draw[->, thick] (s) -- (c);
      \draw[->, thick] (c) -- (e);
      \draw[->, thick] (e) -- (o);
      \node[note, below = 2mm of s, align = center] {trusted:\\ read it};
      \node[note, below = 2mm of c, align = center, text = good]
        {\textbf{checked}\\ by Rocq's kernel};
      \node[note, below = 2mm of e, align = center] {trusted};
      \node[note, below = 2mm of o, align = center] {trusted};
    \end{tikzpicture}
  \end{center}
  \begin{itemize}[<+->]
    \item Who picked the steps does not matter: a person, a script, a model.
    \item A wrong step is \textbf{rejected}, not miscompiled.
  \end{itemize}
  \speakernote{Rocq (formerly Coq) is a proof assistant: a small kernel
    checks every proof. The calculation is a proof, so the optimisation
    steps are outside the trusted base. What remains trusted: that the
    specification says what you mean, the kernel, the extraction to OCaml,
    and the OCaml compiler. [click] Important for later: the steps can come
    from anywhere, including an untrusted search or a learned model. [click]
    If a step is wrong, Rocq refuses it; it cannot silently produce a wrong
    program.}
\end{frame}

\begin{frame}
  \frametitle{Three Obstacles}
  \begin{center}
    \begin{tikzpicture}[alt={Three obstacles side by side: termination,
        because Rocq only accepts recursion it can see terminates; no extra
        axioms, because equality of functions needs one; and clean code,
        because generic definitions extract to code full of unsafe casts.},
      card/.style = {sticky-ish, text width = 33mm, align = center,
        minimum height = 26mm},
      sticky-ish/.style = {draw = golden!80!black, fill = yellow!20,
        rounded corners = 3pt}]
      \node[card] (t) {\textbf{Termination}\\[2mm] \small every function
        must provably stop};
      \node[card, right = 6mm of t] (a) {\textbf{No extra axioms}\\[2mm]
        \small comparing functions needs one};
      \node[card, right = 6mm of a] (c) {\textbf{Clean code}\\[2mm]
        \small generic proofs extract to unsafe casts};
    \end{tikzpicture}
  \end{center}
  \uncover<2->{Haskell has the laws but not termination; Rocq has
    termination but rejects the textbook definitions.}
  \speakernote{Why is this research and not an afternoon of coding?
    Termination: Rocq is a logic, a function that loops would prove
    anything, so it only accepts recursion it can check. Divide and conquer
    recurses on pieces that are not syntactically smaller (partition
    returns new lists). Axioms: the laws are equalities between functions;
    the usual way needs functional extensionality and similar axioms, which
    we avoid so the results stay as trustworthy and as computable as
    possible. Clean code: the textbook encoding of generic data types
    extracts to OCaml full of unsafe casts. [click] Haskell libraries have
    the laws but no termination guarantee.}
\end{frame}

\begin{frame*}
  \frametitle{What the Textbook Encoding Extracts To}
  \badblock{naive-extr.ml}
  \par
  \begin{onlyenv}<2->
    \begin{infobox}
      \code{Obj.magic}: an \textbf{unchecked cast}. Unsafe as soon as this
      code meets other code.
    \end{infobox}
  \end{onlyenv}
  \par
  \speakernote{Real output of Rocq's extraction for the textbook encoding
    of generic data types. OCaml's type system cannot express the types,
    so extraction falls back to Obj.magic, an unchecked cast, everywhere.
    The code is correct as long as it only talks to itself; linked with
    hand-written code it can crash or corrupt memory. Not acceptable for
    verified code meant to be used.}
\end{frame*}

\begin{frame}
  \frametitle{Our Fix: Shapes and Positions}
  \begin{center}
    \begin{tikzpicture}[alt={A data structure seen as a shape with
        positions. The quicksort node Node 3 has two positions, left and
        right, each holding a list. A leaf has no positions. A table on the
        right says which positions are valid in which shape: left and right
        are valid in a node, neither in a leaf.},
      node distance = 6mm and 14mm]
      \node[hot] (s) {shape: $\op{Node}\;3$};
      \node[box, above right = 0mm and 16mm of s] (l) {$[1,2]$};
      \node[box, below right = 0mm and 16mm of s] (r) {$[4,5]$};
      \draw[->] (s) -- node[above, note] {left} (l);
      \draw[->] (s) -- node[below, note] {right} (r);
      \node[hot, below = 14mm of s] (lf) {shape: $\op{Leaf}$};
      \node[note, right = 2mm of lf] {no positions};
      \node[anchor = west] at (7.4, -0.7) {%
        \renewcommand\arraystretch{1.3}%
        \begin{tabular}{lcc}
          \textbf{valid?} & \textbf{left} & \textbf{right} \\
          \code{Node h} & \textcolor{good}{yes} & \textcolor{good}{yes} \\
          \code{Leaf} & \textcolor{bad}{no} & \textcolor{bad}{no} \\
        \end{tabular}};
    \end{tikzpicture}
  \end{center}
  \begin{itemize}[<+->]
    \item One type of positions + a \textbf{yes/no check}: plain OCaml types.
    \item Proofs about the check are \textbf{erased} from the code.
  \end{itemize}
  \speakernote{Containers (Abbott, Altenkirch, Ghani): a data structure is
    a shape plus a function from positions to contents. The textbook
    version gives every shape its own type of positions, which OCaml
    cannot type. [click] We use one type of positions for all shapes, and a
    boolean function saying which positions are valid in a shape. Because
    the check is a boolean, its proofs are unique without axioms
    (Hedberg's theorem). [click] Proofs live in Rocq's erasable universe,
    so extraction drops them: positions become plain values, and the
    contents a plain function. No casts.}
\end{frame}

\begin{frame}
  \frametitle{Termination: Recursion on a Proof}
  \begin{center}
    \begin{tikzpicture}[alt={Left: a call tree of quicksort. Right: the
        termination proof, a finite tree with the same shape. The recursion
        follows the proof. At extraction, the proof is erased and only the
        recursion remains.},
      level distance = 9mm, sibling distance = 18mm,
      every node/.style = {box, font = \small}]
      \begin{scope}
        \node {$x$}
          child {node {$x_1$} child {node {leaf}}}
          child {node {$x_2$} child {node {leaf}} child {node {leaf}}};
        \node[draw = none, fill = none, note] at (0, 0.9) {call tree};
      \end{scope}
      \node[draw = none, fill = none, font = \Large] at (3.6, -0.9)
        {$\leftrightarrow$};
      \begin{scope}[xshift = 72mm, every node/.style = {box, font = \small,
          draw = good, fill = good!8}]
        \node {$\checkmark$}
          child {node {$\checkmark$} child {node {$\checkmark$}}}
          child {node {$\checkmark$} child {node {$\checkmark$}}
            child {node {$\checkmark$}}};
        \node[draw = none, fill = none, note] at (0, 0.9)
          {termination proof};
      \end{scope}
    \end{tikzpicture}
  \end{center}
  \begin{itemize}[<+->]
    \item Proved \textbf{once per split}, e.g.\ ``pieces are shorter''.
    \item The proof is scaffolding: \textbf{erased} at extraction.
  \end{itemize}
  \speakernote{Rocq accepts recursion on a structure that gets smaller. The
    trick: recurse on a proof that the split terminates. The proof is a
    finite tree with the same shape as the call tree, so recursion on it is
    structural. [click] The proof comes from a measure, once per split:
    for quicksort, the pieces are shorter than the input. Every algorithm
    that uses that split gets termination for free. [click] The proof lives
    in the erasable part of Rocq, so the extracted code is just the
    recursion.}
\end{frame}

\begin{frame*}
  \frametitle{No Extra Axioms}
  \outblock{assumptions.txt}
  \begin{itemize}[<+->]
    \item Equality = \emph{same results on equal inputs} (setoids).
    \item Laws become \code{rewrite} steps.
  \end{itemize}
  \speakernote{Rocq can list every assumption a theorem depends on. The
    fusion law and the whole maximum segment sum calculation (later)
    depend on nothing. The quicksort example only depends on Rocq's
    primitive machine integers, which extract to native OCaml integers.
    [click] How we avoid the function extensionality axiom: we never claim
    two functions are identical, only that they give equal results on equal
    inputs, and make every definition respect that. [click] Rocq's
    generalised rewriting then lets us use the laws as rewrite rules.}
\end{frame*}
